\section*{Chapitre \ref{ch:b1:curves} --- Courbes planes}

\begin{solution}{exo:b1:curves:1}
$M(-t) = (t^2, -t^3)$~: réflexion par rapport à l'axe des $x$~; on étudie $t \geq 0$.
$x' = 2t$ et $y' = 3t^2$ sont tous deux $\geq 0$~: la branche va vers
la droite et vers le haut, de $(0,0)$ jusqu'à l'infini. En $t = 0$ la vitesse
s'annule~; les développements $x = t^2$, $y = t^3$ montrent que la tangente est
l'axe des $x$ ($y/x = t \to 0$), avec $y$ changeant de signe tandis que $x \geq 0$~:
un \emph{\omterm{ex:b1:curves:astroid}{point de rebroussement}} pointant vers la gauche. La courbe est la parabole
semi-cubique $y^2 = x^3$.
\end{solution}

\begin{solution}{exo:b1:curves:2}
$M(t + 2\pi) = M(t) + (2\pi, 0)$~: la courbe se répète, translatée d'une
circonférence de roue~; on étudie une période. $x'(t) = 1 - \cos t \geq
0$, $y'(t) = \sin t$~: l'arche monte sur $\intoo{0}{\pi}$, descend sur
$\intoo{\pi}{2\pi}$, culminant en $(\pi, 2)$. Points singuliers
là où $x' = y' = 0$~: $t \in 2\pi\Z$, sur le sol. Au \omterm{def:b1:topology:open}{voisinage} de $t = 0$~:
\[
x(t) = \frac{t^3}{6} + o(t^3),
\qquad
y(t) = \frac{t^2}{2} + o(t^2):
\]
$x$ change de signe, $y \geq 0$, et $\frac{x}{y^{3/2}}$ est borné~: la
tangente est \emph{verticale} (direction $(0,1)$), un \omterm{ex:b1:curves:astroid}{point de rebroussement} où le
point suivi a momentanément une vitesse nulle --- la signature physique
du roulement sans glissement.
\end{solution}

\begin{solution}{exo:b1:curves:3}
En $t = \frac\pi4$~: $M = \bigl(\frac{\sqrt2}{4},
\frac{\sqrt2}{4}\bigr)$, direction tangente $(-\cos t, \sin t) =
\frac{1}{\sqrt2}(-1, 1)$ (\cref{ex:b1:curves:astroid}). Tangente~:
$y - \frac{\sqrt2}{4} = -(x - \frac{\sqrt2}{4})$, c'est-à-dire $x +
y = \frac{\sqrt2}{2}$. Intersections avec les axes~: $\bigl(\frac{\sqrt2}{2},
0\bigr)$ et $\bigl(0, \frac{\sqrt2}{2}\bigr)$~; le segment qui les joint
a pour \omterm{def:b1:curves:arclength}{longueur} $\sqrt{\frac12 + \frac12} = 1$.

(Point général $t$~: la tangente en $(\cos^3 t, \sin^3 t)$ coupe les
axes en $(\cos t, 0)$ et $(0, \sin t)$ --- vérifier que la droite
passant par ces points a pour direction $(-\cos t, \sin t)$ et passe
par $M(t)$ --- et le segment découpé a pour \omterm{def:b1:curves:arclength}{longueur} $\sqrt{\cos^2 t +
\sin^2 t} = 1$.)
\end{solution}

\begin{solution}{exo:b1:curves:4}
$r = \cos 2\theta$~: période $\pi$ en $\theta$, symétrique par rapport aux deux
axes~; $r$ s'annule en $\theta = \pm\frac\pi4$ (tangentes à
l'origine le long des diagonales) et est négatif pour $\theta \in
\intoo{\frac\pi4}{\frac{3\pi}{4}}$, où les points se placent sur le
rayon opposé --- ce qui produit quatre pétales le long des directions des axes
$\theta = 0, \frac\pi2, \pi, \frac{3\pi}{2}$, chacun de rayon
maximal $1$.

$r\cos\theta = 1$ est l'équation $x = 1$~: la \omterm{def:b1:curves:polar}{courbe en polaires} $r =
\frac{1}{\cos\theta}$ est la droite verticale $x = 1$ (décrite une fois pour
$\theta \in \intoo{-\frac\pi2}{\frac\pi2}$).
\end{solution}

\begin{solution}{exo:b1:curves:5}
$r(\theta) = 1 + 2\cos\theta$ s'annule pour $\cos\theta = -\frac12$~:
$\theta = \pm\frac{2\pi}{3}$. Symétrie par rapport à l'axe des $x$~; on étudie
$\theta \in \intcc{0}{\pi}$. Pour $\theta \in
\intco{0}{\frac{2\pi}{3}}$~: $r > 0$, décroissant de $3$ à $0$~:
arc extérieur, entrant à l'origine tangentiellement au rayon $\theta =
\frac{2\pi}{3}$. Pour $\theta \in \intoc{\frac{2\pi}{3}}{\pi}$~: $r <
0$~: les points $r\vec u(\theta)$ se placent sur le rayon opposé (angle
$\theta - \pi \in \intoc{-\frac\pi3}{0}$), avec une distance $\abs r$
croissant de $0$ à $1$~: cela dessine une petite \emph{boucle intérieure}
passant par l'origine. En $\theta = \pi$, $r = -1$ et le point est
$-\vec u(\pi) = (1, 0)$~: la boucle se referme sur le demi-axe positif des $x$.
Tracé~: une grande courbe extérieure en forme de cœur,
de portée maximale $3$ en $\theta = 0$, plus une boucle à l'\omterm{def:b1:topology:closure}{intérieur}
passant par $O$ et $(1,0)$.
\end{solution}

\begin{solution}{exo:b1:curves:6}
Complétons les carrés~:
\[
9(x^2 - 4x) + 25(y^2 - 2y) = 164
\iff 9(x-2)^2 + 25(y-1)^2 = 164 + 36 + 25 = 225 .
\]
En divisant par $225$~: $\frac{(x-2)^2}{25} + \frac{(y-1)^2}{9} = 1$~: une
ellipse de centre $(2, 1)$, de demi-axes $a = 5$ (horizontal) et $b =
3$~; $c = \sqrt{25 - 9} = 4$~: foyers $(2 \pm 4,\, 1) = (-2, 1)$ et
$(6, 1)$~; excentricité $e = \frac45$.
\end{solution}

\begin{solution}{exo:b1:curves:7}
Foyer à l'origine, directrice $D: x = d$ avec $d = \frac pe > 0$.
Pour un point $M = (r\cos\theta, r\sin\theta)$ avec $r > 0$~:
\[
MF = r,
\qquad
d(M, D) = \abs{d - r\cos\theta},
\]
et la condition de \omterm{def:b1:curves:conics}{conique} $MF = e\,d(M, D)$, dans le régime où $M$
est du côté du foyer par rapport à la directrice ($r\cos\theta < d$), s'écrit $r
= e(d - r\cos\theta)$, c'est-à-dire
\[
r(1 + e\cos\theta) = ed = p,
\qquad
r = \frac{p}{1 + e\cos\theta} .
\]
Pour $e < 1$ le dénominateur ne s'annule jamais~: tout $\theta$ est permis
(ellipse). Pour $e = 1$~: $\theta \neq \pi$ (parabole, \omterm{def:b1:topology:open}{ouverte} du côté
opposé à la directrice). Pour $e > 1$~: il faut $1 + e\cos\theta > 0$,
c'est-à-dire $\theta \in \intoo{-\theta_0}{\theta_0}$ avec $\theta_0 =
\arccos\bigl(-\frac1e\bigr)$~: une \emph{branche} de l'hyperbole
(l'autre branche correspond au choix de signe $r < 0$, ou au
second foyer).
\end{solution}

\begin{solution}{exo:b1:curves:8}
\emph{Le jardinier~:} attacher une ficelle de \omterm{def:b1:curves:arclength}{longueur} $2a$ à deux piquets $F$ et
$F'$ et la maintenir tendue avec la pointe traçante $M$~: la contrainte est
exactement $MF + MF' = 2a$, donc la courbe tracée est l'ellipse.

\emph{Propriété de réflexion~:} soit $t \mapsto M(t)$ un paramétrage
régulier et $u(t) = \frac{M(t) - F}{\norm{M(t) - F}}$,
$u'(t)$ le vecteur unitaire analogue vers $F'$. En dérivant
$\norm{M - F} + \norm{M - F'} = 2a$, à l'aide de $\frac{\dd}{\dd
t}\norm{M - F} = \bigl\langle \frac{M - F}{\norm{M-F}},\,
M'\bigr\rangle$ (\omterm{thm:b1:derivative:operations}{règle de la chaîne} sur $\sqrt{\langle\cdot,\cdot\rangle}$)~:
\[
\bigl\langle u + u',\, M'\bigr\rangle = 0 .
\]
Ainsi la direction tangente $M'$ est \omterm{def:b1:euclid:orthogonal}{orthogonale} à la direction $u + u'$
de la bissectrice des deux rayons focaux~: la tangente fait des angles égaux
avec $MF$ et $MF'$. (Un rayon lumineux issu d'un foyer se réfléchit
sur l'ellipse vers l'autre foyer.)
\end{solution}

\begin{solution}{exo:b1:curves:9}
\begin{enumerate}
  \item Domaine~: $\cos 2\theta \geq 0$~: $\theta \in
        \intcc{-\frac\pi4}{\frac\pi4} \cup
        \intcc{\frac{3\pi}{4}}{\frac{5\pi}{4}}$. Symétries~: $\theta
        \mapsto -\theta$ (axe des $x$) et $\theta \mapsto \pi -
        \theta$ (axe des $y$)~: on étudie $\intcc{0}{\frac\pi4}$~; $r$
        décroît de $1$ à $0$. En $\theta = \pm\frac\pi4$~: $r =
        0$, tangentes à l'origine le long des diagonales. La courbe est
        le symbole $\infty$~: deux boucles symétriques se rejoignant en $O$,
        atteignant $(\pm 1, 0)$.
  \item Avec $F, F' = (\pm c, 0)$, $c = \frac{1}{\sqrt 2}$, et $M
        = (r\cos\theta, r\sin\theta)$~:
        \[
        MF^2\, MF'^2
        = \bigl((r\cos\theta - c)^2 + r^2\sin^2\theta\bigr)
        \bigl((r\cos\theta + c)^2 + r^2\sin^2\theta\bigr)
        = (r^2 + c^2)^2 - (2rc\cos\theta)^2 ,
        \]
        par l'identité $(A - B)(A + B) = A^2 - B^2$ avec $A = r^2
        + c^2$, $B = 2rc\cos\theta$. Avec $c^2 = \frac12$~:
        \[
        MF^2 MF'^2 = \Bigl(r^2 + \frac12\Bigr)^2 - 2r^2\cos^2\theta
        = r^4 + r^2\bigl(1 - 2\cos^2\theta\bigr) + \frac14
        = r^4 - r^2\cos 2\theta + \frac14 .
        \]
        Sur la lemniscate, $r^2 = \cos 2\theta$~: les deux premiers termes
        se compensent, laissant $MF^2MF'^2 = \frac14$, c'est-à-dire $MF \cdot MF' =
        \frac12$. Réciproquement, le calcul lu à l'envers montre que
        l'équation du lieu $MF\,MF' = \frac12$ est $r^2 = \cos
        2\theta$ (pour $r \neq 0$~; et $O$ vérifie les deux).
\end{enumerate}
\end{solution}

\begin{solution}{exo:b1:curves:10}
$r = 1 + \cos\theta$, $r' = -\sin\theta$~:
\[
r^2 + r'^2 = 1 + 2\cos\theta + \cos^2\theta + \sin^2\theta
= 2 + 2\cos\theta = 4\cos^2\frac\theta2 ,
\]
donc $\sqrt{r^2 + r'^2} = 2\abs{\cos\frac\theta2}$ et, par la
symétrie par rapport à l'axe des $x$,
\[
L = \int_0^{2\pi} 2\Bigl|\cos\frac\theta2\Bigr|\,\dd\theta
= 2\int_0^{\pi} 2\cos\frac\theta2\,\dd\theta
= 2\Bigl[4\sin\frac\theta2\Bigr]_0^{\pi} = 8 .
\]
Encore un périmètre algébrique, sans $\pi$.
\end{solution}

\begin{solution}{exo:b1:curves:11}
Le point $(a\cos t, b\sin t)$ vérifie l'équation~:
$\cos^2 t + \sin^2 t = 1$. Le vecteur normal de la droite est
$\bigl(\frac{\cos t}a, \frac{\sin t}b\bigr)$, et son produit
avec la vitesse $(-a\sin t, b\cos t)$ vaut $-\sin t\cos t +
\sin t\cos t = 0$~: la droite passe par le point avec la
direction tangente --- c'est la tangente. Pour $(x_0, y_0)$ sur
l'ellipse, écrivons $\cos t = \frac{x_0}a$, $\sin t = \frac{y_0}b$ et
substituons~:
\[
\frac{x\,x_0}{a^2} + \frac{y\,y_0}{b^2} = 1 ,
\]
obtenue à partir de l'équation de l'ellipse en «~dédoublant~» $x^2 \mapsto
x\,x_0$ et $y^2 \mapsto y\,y_0$.
\end{solution}

\begin{solution}{exo:b1:curves:12}
\begin{enumerate}
  \item $r' = k\eu^{k\theta}$, donc $\tan V = \frac{r}{r'} =
        \frac1k$~: constant. La spirale coupe tout rayon issu de
        l'origine sous le même angle $V = \arctan\frac1k$.
  \item En notation complexe la spirale est $\{\eu^{k\theta}
        \eu^{\iu\theta} : \theta \in \R\}$. La rotation d'angle $c$
        multiplie par $\eu^{\iu c}$~:
        \[
        \eu^{k\theta}\eu^{\iu(\theta + c)}
        = \eu^{-kc}\,\eu^{k(\theta + c)}\eu^{\iu(\theta+c)} ,
        \]
        et lorsque $\theta + c$ parcourt $\R$ ceci décrit
        $\eu^{-kc}$ fois la spirale~: rotation $=$ homothétie. Aucune
        autre courbe lisse que les droites et les cercles n'a cette
        propriété.
  \item $\sqrt{r^2 + r'^2} = \sqrt{1 + k^2}\,\eu^{k\theta}$, donc
        sur $\intcc{A}{\theta_0}$ la \omterm{def:b1:curves:arclength}{longueur} vaut
        $\frac{\sqrt{1+k^2}}{k}\bigl(\eu^{k\theta_0} -
        \eu^{kA}\bigr)$, et quand $A \to -\infty$~:
        \[
        L = \frac{\sqrt{1 + k^2}}{k}\,\eu^{k\theta_0} < \infty :
        \]
        une infinité de tours autour de l'origine, et une \omterm{def:b1:curves:arclength}{longueur} totale
        finie (les tours décroissent géométriquement).
\end{enumerate}
\end{solution}

\begin{solution}{pb:b1:curves:1}
\textbf{1.} Rouler sans glisser signifie que le point de contact a
parcouru une distance égale à l'arc de roue déroulé~: après
avoir tourné de $t$, le centre est en $(Rt, R)$. Le point marqué se trouve
sur la jante, à l'angle $t$ \emph{en arrière} de la verticale descendante
(la roue tourne dans le sens horaire en avançant)~:
\[
M(t) = \Omega(t) + R(-\sin t, -\cos t)
= \bigl(R(t - \sin t),\ R(1 - \cos t)\bigr),
\]
ce qui est correct en $t = 0$ ($M = (0,0)$) et en $t = \pi$ ($M =
(\pi R, 2R)$, le sommet).

\textbf{2.} $f'(t) = R(1 - \cos t,\ \sin t)$ et
\[
\norm{f'}^2 = R^2\bigl((1-\cos t)^2 + \sin^2 t\bigr)
= 2R^2(1 - \cos t) = 4R^2\sin^2\frac t2 ,
\]
donc $\norm{f'} = 2R\abs{\sin\frac t2}$. Points singuliers en $t \in
2\pi\Z$~: les \omterm{ex:b1:curves:astroid}{points de rebroussement} sur le sol (\cref{exo:b1:curves:2})~; le
sommet de l'arche est en $t = \pi$, où la vitesse $2R$ est maximale.

\textbf{3.} Factorisations par l'angle moitié~:
\[
f'(t) = 2R\sin\frac t2\,\Bigl(\sin\frac t2,\ \cos\frac t2\Bigr),
\quad
M - C = 2R\sin\frac t2\,\Bigl(-\cos\frac t2,\ \sin\frac t2\Bigr),
\]
\[
T - M = \bigl(R\sin t,\ R(1 + \cos t)\bigr)
= 2R\cos\frac t2\,\Bigl(\sin\frac t2,\ \cos\frac t2\Bigr).
\]
Pour $0 < t < 2\pi$, $\sin\frac t2 \neq 0$~: $M - C$ est \omterm{def:b1:euclid:orthogonal}{orthogonal}
à la direction tangente $\bigl(\sin\frac t2, \cos\frac
t2\bigr)$ (leur produit vaut $-\sin\frac t2\cos\frac t2 +
\sin\frac t2\cos\frac t2 = 0$), et $T - M$ lui est parallèle.
La corde vers le sommet de la roue \emph{est} la tangente~; la
corde vers le point de contact est la normale.

\textbf{4.} À chaque instant la roue pivote autour de son point de
contact (ce point a une vitesse nulle~: roulement sans glissement),
donc tout point solidaire de la roue se déplace orthogonalement à la droite
qui le joint à $C$, avec une vitesse (pour une vitesse angulaire $1$) égale à
cette distance. Vérification~: $\norm{M - C} = 2R\abs{\sin\frac t2} =
\norm{f'(t)}$.

\textbf{5.} $2R\,y(t) = 2R\cdot 2R\sin^2\frac t2 =
4R^2\sin^2\frac t2 = \norm{f'(t)}^2$. La vitesse à la hauteur $y$ est
$\sqrt{2R\,y}$ --- formellement la loi $v = \sqrt{2gh}$ de la chute
\omterm{def:b1:vspaces:free}{libre}, le sol jouant le rôle du plafond~; la partie IV transforme cette
observation en horlogerie.

\textbf{6.} D'après la \cref{def:b1:curves:arclength} et la question 2
($\sin\frac t2 \geq 0$ sur $\intcc{0}{2\pi}$)~:
\[
L = \int_0^{2\pi} 2R\sin\frac t2\,\dd t
= 2R\Bigl[-2\cos\frac t2\Bigr]_0^{2\pi} = 2R(2 + 2) = 8R .
\]
Théorème de Wren~: exactement quatre diamètres.

\textbf{7.} $s(t) = \int_0^t 2R\sin\frac u2\,\dd u =
4R\bigl(1 - \cos\frac t2\bigr)$~; $s(2\pi) = 4R(1+1) = 8R$,
ce qui est cohérent.

\textbf{8.} $\sigma(t) = \abs{s(t) - s(\pi)} = \abs{4R(1 -
\cos\frac t2) - 4R} = 4R\abs{\cos\frac t2}$, donc $\sigma^2 =
16R^2\cos^2\frac t2$~; et $2R - y = R(1 + \cos t) =
2R\cos^2\frac t2$, d'où $8R(2R - y) = 16R^2\cos^2\frac t2 =
\sigma^2$.

\textbf{9.} Au \omterm{ex:b1:curves:astroid}{point de rebroussement} $t = 0$ (ou $2\pi$)~: $\sigma = 4R$, la moitié
de $8R$~: le sommet coupe l'arche en deux. Dans les deux calculs, la
vitesse est $\abs{\text{polynôme trigonométrique en } t/2}$, dont
une primitive est encore trigonométrique~: la \omterm{def:b1:curves:arclength}{longueur} est une
différence de \emph{valeurs} de cosinus --- des nombres rationnels
multipliés par $R$ --- sans aucun arc de cercle à mesurer, d'où l'absence de $\pi$.
Le cercle, lui, a une vitesse constante, si bien que son \omterm{thm:b1:integration:def}{intégrale} de \omterm{def:b1:curves:arclength}{longueur}
produit la \omterm{def:b1:curves:arclength}{longueur} totale de l'\omterm{prop:b1:reals:intervals}{intervalle} $2\pi$~; la vitesse de la cycloïde
s'annule aux extrémités et s'intègre algébriquement.

\textbf{10.} L'arche est décrite avec $x$ croissant de $0$ à
$2\pi R$ ($x' = R(1 - \cos t) \geq 0$, ne s'annulant qu'en des
points isolés). La substitution $x = x(t)$ dans l'\omterm{thm:b1:integration:def}{intégrale} d'aire
$\int_0^{2\pi R} y\,\dd x$ (\cref{thm:b1:integration:parts})
donne $A = \int_0^{2\pi} y(t)\,x'(t)\,\dd t$.

\textbf{11.}
\[
A = \int_0^{2\pi} R(1 - \cos t)\cdot R(1 - \cos t)\,\dd t
= R^2\int_0^{2\pi}\bigl(1 - 2\cos t + \cos^2 t\bigr)\dd t
= R^2\Bigl(2\pi - 0 + \pi\Bigr) = 3\pi R^2 ,
\]
en utilisant $\int_0^{2\pi}\cos^2 = \pi$. Exactement trois aires de roue~:
la balance de Galilée disait «~environ $3$~»~; le calcul de Roberval
dit «~exactement~».

\textbf{12.} Avec $x = \cos^3 t$, $y = \sin^3 t$~: $y\,x' =
-3\sin^4 t\cos^2 t$. Linéarisons~:
\[
\sin^4 t\cos^2 t = \frac{1 - \cos 2t}{2}\cdot\frac{\sin^2
2t}{4} = \frac{\sin^2 2t}{8} - \frac{\sin^2 2t\cos 2t}{8} ,
\]
et sur $\intcc{0}{2\pi}$~: $\int\sin^2 2t = \pi$, $\int \sin^2
2t\cos 2t = \bigl[\frac{\sin^3 2t}{6}\bigr] = 0$. Donc $\oint
y\,\dd x = -3\cdot\frac\pi8$, et l'aire enfermée vaut
$\frac{3\pi}8$ (le signe enregistre l'orientation dans le sens
trigonométrique).

\textbf{13.} Pour $(\cos t, \sin t)$ avec $t$ allant de $\pi$ à $0$,
$x$ croît de $-1$ à $1$ et
\[
\int_\pi^0 \sin t\cdot(-\sin t)\,\dd t
= \int_0^\pi \sin^2 t\,\dd t = \frac\pi2 :
\]
la formule rend l'aire du demi-disque supérieur, comme il se doit.

\textbf{14.} $x' = R(1 + \cos t) = 2R\cos^2\frac t2$, $y' =
R\sin t = 2R\sin\frac t2\cos\frac t2$, donc $\norm{f'} =
2R\cos\frac t2$ (positif sur $\intcc{-\pi}{\pi}$). Arc depuis
le fond~: $s(t) = \int_0^t 2R\cos\frac u2\,\dd u =
4R\sin\frac t2$. Alors
\[
y = R(1 - \cos t) = 2R\sin^2\frac t2
= 2R\Bigl(\frac{s}{4R}\Bigr)^{\!2} = \frac{s^2}{8R} .
\]

\textbf{15.} Conservation de l'énergie avec $v = \frac{\dd s}{\dd\tau}$
et $y = \frac{s^2}{8R}$, $y_0 = \frac{s_0^2}{8R}$~:
\[
\Bigl(\frac{\dd s}{\dd\tau}\Bigr)^{\!2} = 2g(y_0 - y)
= \frac{2g}{8R}\bigl(s_0^2 - s^2\bigr)
= \frac{g}{4R}\bigl(s_0^2 - s^2\bigr).
\]

\textbf{16.} En dérivant en $\tau$~: $2s's'' =
-\frac{g}{4R}\,2ss'$, donc partout où $s' \neq 0$ (donc partout
par \omterm{def:b1:continuity:continuous}{continuité}), $s'' = -\frac{g}{4R}s$~: l'oscillateur harmonique.
D'après le \cref{thm:b1:diffeq:homogeneous2}, $s(\tau) =
A\cos\omega\tau + B\sin\omega\tau$ avec $\omega =
\sqrt{g/(4R)}$~; les conditions initiales $s(0) = s_0$, $s'(0) = 0$
donnent $s(\tau) = s_0\cos\omega\tau$.

\textbf{17.} La bille atteint le fond quand $s = 0$, c'est-à-dire en
$\omega\tau = \frac\pi2$~:
\[
T_{\downarrow} = \frac{\pi}{2\omega}
= \frac\pi2\sqrt{\frac{4R}{g}} = \pi\sqrt{\frac Rg}\,,
\]
indépendant de $s_0$~: lâchées de n'importe où dans la cuvette, toutes
les billes arrivent au même instant --- la tautochrone.

\textbf{18.} En substituant $s = s_0\cos\omega\tau$~: le membre de gauche
vaut $s_0^2\omega^2\sin^2\omega\tau$ et le membre de droite vaut
$\frac{g}{4R}s_0^2(1 - \cos^2\omega\tau) =
s_0^2\omega^2\sin^2\omega\tau$~: égalité exacte, donc aucune
solution parasite n'a été introduite. Pour un arc de cercle, $y$ n'est pas
proportionnel à $s^2$ ($y = \ell(1 - \cos\frac s\ell) =
\frac{s^2}{2\ell} - \frac{s^4}{24\ell^3} + \dots$)~: le
terme de rappel n'est qu'\emph{approximativement} \omterm{def:b1:linmaps:def}{linéaire}, si bien que la
période d'un pendule circulaire dérive avec l'amplitude, alors que celle de la
cycloïde est rigoureusement constante.

\textbf{19.} $\pi\sqrt{R/g} = 1$ donne $R = \frac{g}{\pi^2}
\approx \frac{9.81}{9.87} \approx 0.994$ m --- presque exactement
un mètre. L'oscillation complète prend $\frac{2\pi}\omega =
2\pi\sqrt{4R/g}$, qui est précisément la formule des petits angles
$2\pi\sqrt{\ell/g}$ pour un pendule de \omterm{def:b1:curves:arclength}{longueur} $\ell = 4R$~:
Huygens suspendait son pendule entre des joues cycloïdales exactement
de cette proportion.

\textbf{20.} En $t = \frac\pi2$, $R = 1$~: $M = \bigl(\frac\pi2 -
1,\ 1\bigr)$, $T = \bigl(\frac\pi2,\ 2\bigr)$, donc $T - M = (1,
1)$. Et $f'(\frac\pi2) = (1 - 0,\ 1) = (1, 1)$~: la corde vers
le sommet est la vitesse, comme promis.

\textbf{21.} Pour la trochoïde, $x'(t) = R - d\cos t$ et $y'(t)
= d\sin t$. Si $d < R$~: $x' \geq R - d > 0$, le point avance
toujours~; la vitesse ne s'annule jamais (sa première composante est
positive)~: aucun point singulier, une ondulation régulière. Si $d > R$~: $x'(0)
= R - d < 0 < R + d = x'(\pi)$, donc le point recule
au \omterm{def:b1:topology:open}{voisinage} des instants de contact et avance ailleurs~: la courbe
se recoupe en formant des boucles. Un point du boudin d'une roue de
chemin de fer (sous le champignon du rail, $d > R$) recule à chaque
tour.

\textbf{22.} La corde du \omterm{ex:b1:curves:astroid}{point de rebroussement} $(\pi R, 2R)$ à l'origine
a pour \omterm{def:b1:curves:arclength}{longueur} $L = R\sqrt{\pi^2 + 4}$ et fait un angle $\alpha$ avec
l'horizontale, avec $\sin\alpha = \frac{2R}{L}$. En glissant depuis le repos avec une
accélération constante $g\sin\alpha$~: $L = \frac12 g\sin\alpha\,T^2$, donc
\[
T = \sqrt{\frac{2L}{g\sin\alpha}} = \sqrt{\frac{2L^2}{2Rg}}
= \frac{L}{\sqrt{Rg}} = \sqrt{\pi^2 + 4}\,\sqrt{\frac Rg}
\approx 3.72\sqrt{\frac Rg}\,,
\]
contre $\pi\sqrt{R/g} \approx 3.14\sqrt{R/g}$ pour la cycloïde~:
le chemin courbe est plus rapide. Il est en fait le plus rapide possible
--- la brachistochrone --- un théorème du calcul des
variations, hors de portée de ce volume.

\textbf{23.} $\dfrac{\dd y}{\dd x} = \dfrac{y'}{x'} =
\dfrac{\sin t}{1 - \cos t} = \cot\dfrac t2$ (formules de l'angle
moitié). Alors
\[
\frac{\dd^2y}{\dd x^2}
= \frac{\dd}{\dd t}\Bigl(\cot\frac t2\Bigr)\cdot\frac{1}{x'(t)}
= -\frac{1}{2\sin^2\frac t2}\cdot\frac{1}{2R\sin^2\frac t2}
= -\frac{1}{4R\sin^4\frac t2} < 0 :
\]
concave sur toute l'arche~; quand $t \to 0^+$ ou $2\pi^-$, la
pente $\cot\frac t2 \to \pm\infty$~: tangentes verticales aux
\omterm{ex:b1:curves:astroid}{points de rebroussement}.

\textbf{24.} La direction de la corde est $T - M \parallel
\bigl(\sin\frac t2, \cos\frac t2\bigr)$, qui tend vers $(0, 1)$
quand $t \to 0^{+}$~: la tangente au \omterm{ex:b1:curves:astroid}{point de rebroussement} est verticale ---
retrouvée par la géométrie pure, sans aucun développement de Taylor.

\textbf{25.} (i) Le théorème de Wren, $L = 8R$~; l'aire de Roberval, $A =
3\pi R^2$ (le rapport $3$ conjecturé par Galilée)~; la tautochrone de
Huygens, $T_\downarrow = \pi\sqrt{R/g}$. (ii) Tout le
calcul a reposé sur les factorisations par l'angle moitié $f' =
2R\sin\frac t2(\sin\frac t2, \cos\frac t2)$ et leur analogue dans la
cuvette --- une seule identité alimentant la tangente, la \omterm{def:b1:curves:arclength}{longueur} et l'horloge.
(iii) La relation intrinsèque $y = s^2/8R$ a converti l'équation
d'énergie en $s'' = -\frac{g}{4R}s$, une équation \omterm{def:b1:linmaps:def}{linéaire}
à coefficients constants dont les solutions sont exactement
isochrones. (iv) La partie I a utilisé le calcul différentiel du
\cref{ch:b1:derivative}, les parties II et III l'\omterm{thm:b1:integration:def}{intégrale} du
\cref{ch:b1:integration}, la partie IV les équations différentielles du
\cref{ch:b1:diffeq} --- la cycloïde est le programme de ce livre
enroulé en une seule courbe.
\end{solution}
